% ==============================================================================
% The Grand Daily Tournament — 2026-09-11 Match (Week 2 Day 5)
% Locked template: EB Garamond + Garamond-Math + STKaiti (v5.0 - 6 Pages)
% Feature: Full 6-Page Functional Layout with Dedicated Calculation Topic
% ==============================================================================
\documentclass[a4paper,11pt]{article}

% ——— Engine & Font Setup ———
\usepackage{fontspec}
\usepackage{amsmath}
\usepackage{unicode-math}
\defaultfontfeatures{Ligatures=TeX}
\setmainfont{EBGaramond}[
  Path           = {c:/texlive/2024/texmf-dist/fonts/opentype/public/ebgaramond/},
  Extension      = .otf,
  UprightFont    = *-Regular,
  ItalicFont     = *-Italic,
  BoldFont       = *-SemiBold,
  BoldItalicFont = *-SemiBoldItalic,
  Ligatures      = TeX,
  Numbers        = OldStyle
]
\setsansfont{LibertinusSans}[
  Path        = {c:/texlive/2024/texmf-dist/fonts/opentype/public/libertinus-fonts/},
  Extension   = .otf,
  UprightFont = *-Regular,
  ItalicFont  = *-Italic,
  BoldFont    = *-Bold,
  Scale       = MatchLowercase
]
\setmathfont{Garamond-Math.otf}[
  Path  = {c:/texlive/2024/texmf-dist/fonts/opentype/public/garamond-math/},
  Scale = MatchLowercase
]

\usepackage{xeCJK}
\xeCJKsetup{
  CJKmath      = false,
  PunctStyle   = kaiming,
  CheckSingle  = true,
  AutoFakeBold = true,
  AutoFakeSlant= false
}
\setCJKmainfont{STKaiti}[
  Path        = {C:/Windows/Fonts/},
  UprightFont = STKAITI.TTF,
  BoldFont    = STKAITI.TTF,
  ItalicFont  = STKAITI.TTF,
  Script      = Default
]
\setCJKsansfont{Microsoft YaHei}
\setCJKmonofont{FangSong}

\usepackage[margin=1.5cm,headheight=14pt,headsep=10pt,footskip=18pt]{geometry}
\usepackage{booktabs,enumitem,xcolor,graphicx,tabularx}

\usepackage[breakable,skins,hooks]{tcolorbox}
\usepackage{tikz}
\usetikzlibrary{calc,arrows.meta}
\usepackage{fancyhdr,lastpage}
\usepackage{microtype}

% ——— TeXtured Colour Palette ———
\definecolor{navy}{HTML}{0F172A}
\definecolor{slate}{HTML}{1E293B}
\definecolor{royal}{HTML}{1E40AF}
\definecolor{mist}{HTML}{64748B}
\definecolor{border}{HTML}{E2E8F0}
\definecolor{paper}{HTML}{F8FAFC}
\definecolor{forest}{HTML}{065F46}
\definecolor{amber}{HTML}{D97706}
\definecolor{wine}{HTML}{991B1B}
\definecolor{ice}{HTML}{EFF6FF}

% ——— Typography Tuning ———
\linespread{1.10}
\setlength{\parindent}{0pt}
\setlength{\parskip}{3.5pt plus 1pt minus 1pt}
\setlength{\abovedisplayskip}{5pt plus 2pt minus 2pt}
\setlength{\belowdisplayskip}{5pt plus 2pt minus 2pt}

% ——— Header / Footer ———
\pagestyle{fancy}\fancyhf{}
\fancyhead[L]{\textcolor{mist}{\footnotesize\sffamily 2027 Theoretical Physics · 604 Calculus \& 811 Quantum Mechanics}}
\fancyhead[R]{\textcolor{slate}{\footnotesize\sffamily\bfseries Daily Tournament · 2026-09-11}}
\fancyfoot[L]{\textcolor{mist}{\footnotesize\itshape ``Calculate until the very last step.''}}
\fancyfoot[R]{\textcolor{mist}{\footnotesize\sffamily Page \thepage\ of \pageref{LastPage}}}
\renewcommand{\headrulewidth}{0.4pt}
\renewcommand{\headrule}{\color{border}\hrule width\headwidth height\headrulewidth \vskip-\headrulewidth}

% ——— tcolorbox Environments ———
\newtcolorbox{briefing}[1]{%
  enhanced, blanker, breakable,
  left=4mm, right=3mm, top=3mm, bottom=3mm,
  borderline west={2.5pt}{0pt}{royal},
  colback=paper,
  fonttitle=\sffamily\bfseries\color{slate},
  title={#1},
  attach title to upper={\par\vspace{2mm}}
}

\newtcolorbox{problem}[2]{%
  enhanced, breakable,
  colback=white, colframe=border, boxrule=0.6pt, arc=2mm,
  left=4mm, right=4mm, top=2.8mm, bottom=2.8mm,
  fonttitle=\sffamily\bfseries\color{slate},
  title={\raisebox{0pt}{\color{royal}\rule{3pt}{10pt}}\hspace{4pt}#1\hfill{\normalfont\footnotesize\color{mist}\itshape #2}},
  colbacktitle=paper, coltitle=slate,
  titlerule=0pt, toptitle=2mm, bottomtitle=2mm
}

\newtcolorbox{solution}[1]{%
  enhanced, breakable,
  colback=white, colframe=border, boxrule=0.6pt, arc=2mm,
  left=4mm, right=4mm, top=2.8mm, bottom=2.8mm,
  fonttitle=\sffamily\bfseries\color{forest},
  title={\raisebox{0pt}{\color{forest}\rule{3pt}{10pt}}\hspace{4pt}#1},
  colbacktitle=paper, coltitle=forest,
  titlerule=0pt, toptitle=2mm, bottomtitle=2mm
}

\newtcolorbox{debrief}[1]{%
  enhanced, breakable,
  colback=white, colframe=border, boxrule=0.6pt, arc=2mm,
  left=3.5mm, right=3.5mm, top=2.5mm, bottom=2.5mm,
  fonttitle=\sffamily\bfseries\color{slate},
  title={\raisebox{0pt}{\color{amber}\rule{3pt}{10pt}}\hspace{4pt}#1},
  colbacktitle=paper, coltitle=slate,
  titlerule=0pt, toptitle=2mm, bottomtitle=2mm
}

\newtcolorbox{notebox}{%
  enhanced, blanker,
  left=3mm, right=2mm, top=1.5mm, bottom=1.5mm,
  borderline west={2pt}{0pt}{royal!60},
  colback=ice, fontupper=\small
}

\newtcolorbox{warnbox}{%
  enhanced, blanker,
  left=3mm, right=2mm, top=1.5mm, bottom=1.5mm,
  borderline west={2pt}{0pt}{wine},
  colback=wine!4, fontupper=\small
}

\newcommand{\checkbox}{\ensuremath{\mdlgwhtsquare}}
\newcommand{\alerttag}[1]{\textcolor{wine}{\sffamily\bfseries [#1]}}

% ==============================================================================
\begin{document}

% ==============================================================================
% PAGE 1: COVER
% ==============================================================================
\begin{titlepage}
\thispagestyle{empty}
\begin{tikzpicture}[remember picture,overlay]
  \fill[paper] (current page.south west) rectangle (current page.north east);

  \fill[navy] (current page.north west) rectangle ($(current page.north east)+(0,-3.4cm)$);
  \fill[royal] ($(current page.north west)+(0,-3.4cm)$) rectangle ($(current page.north east)+(0,-3.55cm)$);
  \fill[amber] ($(current page.north west)+(0,-3.55cm)$) rectangle ($(current page.north east)+(0,-3.62cm)$);

  \node[anchor=west, text=white] at ($(current page.north west)+(2.2cm,-1.5cm)$) {
    {\fontsize{12}{14}\selectfont\sffamily\bfseries THE GRAND DAILY TOURNAMENT \textperiodcentered\ 2027}
  };
  \node[anchor=west, text=white!50] at ($(current page.north west)+(2.2cm,-2.3cm)$) {
    {\fontsize{8.5}{11}\selectfont\sffamily UCAS \textperiodcentered\ HIAIS \textperiodcentered\ THEORETICAL PHYSICS ADVANCED TRAINING DOSSIER}
  };

  \begin{scope}[shift={($(current page.center)+(0,1.2cm)$)}]
    \draw[line width=0.5pt, border!80, dashed] (0,0) circle (3.3cm);
    \draw[line width=0.7pt, border] (0,0) circle (2.5cm);
    \draw[line width=0.4pt, royal!30] (0,0) circle (1.6cm);
    \draw[line width=1.2pt, slate!60] (-3.8,-2.2) -- (3.8,2.2);
    \draw[line width=1.2pt, slate!60] (-3.8,2.2) -- (3.8,-2.2);
    \draw[line width=0.7pt, royal!65, -{Stealth[length=2.5mm]}] (-4.2,0) -- (4.2,0)
      node[right, font=\footnotesize\sffamily\color{mist}] {$x,\, r$};
    \draw[line width=0.7pt, royal!65, -{Stealth[length=2.5mm]}] (0,-3.6) -- (0,3.6)
      node[above, font=\footnotesize\sffamily\color{mist}] {$\hat{H},\, E_n$};
    \draw[line width=1.3pt, royal, smooth, samples=120, domain=-3.14:3.14]
      plot ({\x}, {1.1*sin(2.5*\x r)*exp(-\x*\x/5.5)});
    \draw[line width=1.1pt, amber!85!black, smooth, samples=120, domain=-3.14:3.14]
      plot ({\x}, {-0.9*cos(2.5*\x r)*exp(-\x*\x/5.5)});
    \fill[navy] (0,0) circle (3.5pt);
    \fill[amber] (0,0) circle (2pt);
    \fill[royal] (2.12, 1.13) circle (2.2pt);
    \fill[wine] (-1.8, -1.05) circle (2.2pt);
    \fill[forest] (1.0, -1.35) circle (2.2pt);
  \end{scope}

  \node[anchor=center] at ($(current page.center)+(0,-3.4cm)$) {
    \begin{minipage}{0.82\textwidth}
      \centering
      {\fontsize{32}{38}\selectfont\bfseries\color{navy} Daily Tournament}\\[5pt]
      {\fontsize{12}{15}\selectfont\sffamily\bfseries\color{royal}%
        GRAND ARENA OF THEORETICAL PHYSICS}\\[12pt]
      {\color{navy}\rule{0.5\textwidth}{1.2pt}}\\[10pt]
      {\fontsize{10.5}{14}\selectfont\color{slate}\itshape
        ``Today's contestants: 604 Calculus Symmetries $\times$ 811 Perturbation Algebra''\\[3pt]
        {\small\color{mist}\upshape Closed-book \textperiodcentered\
        Pen-and-paper derivation \textperiodcentered\ No shortcuts}}
    \end{minipage}
  };

  \node[anchor=south] at ($(current page.south)+(0,1.6cm)$) {
    \begin{minipage}{0.84\textwidth}
      \begin{tcolorbox}[
        enhanced, colback=white, colframe=border, boxrule=0.6pt, arc=2mm,
        left=4mm, right=4mm, top=3.5mm, bottom=3.5mm
      ]
        \renewcommand{\arraystretch}{1.25}
        \sffamily\small
        \begin{tabularx}{\linewidth}{@{}p{2.6cm}X@{}}
          \textbf{Date} & 2026-09-11 \quad\textbar\quad Week 2 (Day 5) \quad\textbar\quad Level 4: Angular Momentum Master\\
          \textbf{Modules} & \textbf{604}: Double Integral Symmetries \& Green Singularities \quad\textbar\quad \textbf{811}: Perturbation Theory \& 3D Oscillator\\
          \textbf{Error Protocol} & 5-Factor: Knowledge\,\textbar\, Modelling\,\textbar\, Logic\,\textbar\, Calculation\,\textbar\, Time\\
          \textbf{Target} & Full 6-Page Combat Edition \quad\textbar\quad Standard 80 min Complete Run\\
        \end{tabularx}
      \end{tcolorbox}
    \end{minipage}
  };
\end{tikzpicture}
\end{titlepage}

% ==============================================================================
% PAGE 2: §1 TACTICAL BRIEFING & WARM-UP RETRIEVAL
% ==============================================================================
\clearpage

\begin{center}
  {\LARGE\bfseries\color{navy} Theoretical Physics Training \textperiodcentered\ Daily Tournament}\\[3pt]
  {\small\sffamily\color{mist} 2026-09-11 \quad\textbar\quad Week 2 (Day 5) \quad\textbar\quad Phase: Multivariable Field Integrals \& Perturbation Algebra}\\[-4pt]
  {\color{navy}\rule{\textwidth}{0.8pt}}\\[-11pt]
  {\color{border}\rule{\textwidth}{0.3pt}}
\end{center}

\vspace{-6pt}

\begin{briefing}{\S 1\quad Tactical Briefing \& Warm-Up Retrieval}
\begin{itemize}[leftmargin=1.5em, itemsep=2pt, topsep=1pt]
  \item \textbf{604 今日靶向}：熟练应用二重积分普通轴对称性（奇零偶倍）与轮换对称性（对调相加折半心法）；掌握极坐标面积微元 $dxdy = rdrd\theta$ 的行列式雅可比来源；透彻理解格林公式的环量旋度本质，熟练应用非闭合曲线补线法与奇点挖洞法。
  \item \textbf{811 今日靶向}：掌握非简并微扰论核心推导逻辑（左乘投影消项）；掌握一级能量期望值 $E_n^{(1)} = \langle n^{(0)}|H'|n^{(0)}\rangle$、二级能量修正公式与能级互斥图像（基态二阶恒负）；运用三维谐振子升降算符坐标代数秒杀 2015 压轴真题中间态矩阵元。
  \item \textbf{Spaced Retrieval 间隔提取 (Day +3 滚动回溯)}：\textbf{604 (Day 3 回溯)} 泰勒展开交叉项定阶原则与余项严控；\textbf{811 (Day 3 回溯)} 氢原子 $n^2$ 简并度几何求和与 Runge-Lenz 动力学对称性。
  \item \textbf{考试规则}：闭卷、独立手推、严禁查阅笔记；先在下方草稿区完成 5--8 分钟白纸公式破冰与间隔提取。
\end{itemize}

\vspace{1mm}
\textbf{Warm-Up Retrieval Checklist} (5--8\,min): without notes, write on scratch paper:\\
\checkbox\ \alerttag{604 Polar} Jacobian: $J = \frac{\partial(x,y)}{\partial(r,\theta)} = r \implies dxdy = rdrd\theta$ \quad
\checkbox\ \alerttag{604 Green} Singularity loop: $\oint_L \frac{-ydx+xdy}{x^2+y^2} = 2\pi$ (counter-clockwise) \quad
\checkbox\ \alerttag{811 Perturb} 1st order energy: $E_n^{(1)} = \langle n^{(0)}|H'|n^{(0)}\rangle$ \quad
\checkbox\ \alerttag{811 Perturb} 2nd order energy: $E_n^{(2)} = \sum_{m \ne n} \frac{|\langle m^{(0)}|H'|n^{(0)}\rangle|^2}{E_n^{(0)}-E_m^{(0)}} \le 0$ for ground state \quad
\checkbox\ \alerttag{811 Ladder} Position operator: $\hat{x} = \sqrt{\frac{\hbar}{2m\omega}}(a + a^\dagger)$, Selection rule $\Delta n = \pm 1$
\end{briefing}

\vspace{4pt}

\begin{tcolorbox}[
  enhanced, colback=white, colframe=border, boxrule=0.6pt, arc=2mm,
  left=4mm, right=4mm, top=3mm, bottom=3mm,
  fonttitle=\sffamily\bfseries\color{slate},
  title={\raisebox{0pt}{\color{royal}\rule{3pt}{10pt}}\hspace{4pt}Warm-Up Retrieval Workspace \textbar\ 白纸破冰与间隔提取草稿区},
  colbacktitle=paper, coltitle=slate,
  titlerule=0pt, toptitle=2mm, bottomtitle=2mm
]
\textcolor{mist}{\footnotesize\itshape 5--8 分钟闭卷盲写：默写极坐标雅可比微元推导、格林公式奇点积分，写出非简并一阶/二阶微扰公式及升降算符坐标表达式。}
\vspace{9.0cm}
\end{tcolorbox}

% ==============================================================================
% PAGE 3: §2 CORE PROBLEM SET
% ==============================================================================
\clearpage

{\Large\sffamily\bfseries\color{navy} \S 2\quad Core Problem Set}\hfill{\small\sffamily\color{mist}\itshape Complete independently under timed conditions}

\vspace{8pt}

\begin{problem}{Problem M1 \textbar\ 604 Calculus: Multivariable Integral \& Symmetry}{Suggested time: 25\,min}
设闭区域 $D = \{(x,y) \mid x^2 + y^2 \le 1, \, x \ge 0, \, y \ge 0\}$ 为第一象限的四分之一单位圆盘。
试计算二重积分：
\[
  I = \iint_D \frac{x^3 + x^2 y}{x^2 + y^2}\,dxdy
\]
\textit{提示：注意观察区域 $D$ 的对称轴与被积函数项的结构，合理拆项与转化极坐标系，熟练应用 Wallis 公式与代数化简。}
\end{problem}

\vspace{6pt}

\begin{problem}{Problem Q1 \textbar\ 811 Quantum Mechanics: 3D Oscillator Perturbation (2015 UCAS)}{Suggested time: 35\,min}
质量为 $m$ 的三维各向同性谐振子，未受微扰时的哈密顿量为：
\[
  H_0 = \sum_{i=x,y,z} \left( \frac{p_i^2}{2m} + \frac{1}{2}m\omega^2 x_i^2 \right)
\]
未微扰基态记为 $|000\rangle = |0\rangle_x |0\rangle_y |0\rangle_z$，对应未微扰基态能量为 $E_{000}^{(0)} = \frac{3}{2}\hbar\omega$。
现体系受到微弱的立方微扰作用：
\[
  H' = \lambda (x^2 y + x y^2)
\]
其中 $\lambda$ 为微小的实耦合常数。
\begin{enumerate}[leftmargin=1.8em, itemsep=2pt, topsep=1pt]
  \item 试从算符宇称（奇偶性）或升降算符代数性质，严格证明基态能量的一级修正 $E_0^{(1)} = 0$；
  \item 试利用定态非简并微扰理论，精确计算基态能量的二级修正 $E_0^{(2)}$；
  \item 根据计算结果，讨论基态能量在微扰开启后发生上移还是下移，并从微扰论能级互斥图像阐述其物理机制。
\end{enumerate}
\end{problem}

\vspace{6pt}

\begin{problem}{Extension \textbar\ Conceptual Analysis \& Spaced Retrieval}{Optional \textperiodcentered\ Suggested time: 10\,min}
\begin{enumerate}[leftmargin=1.8em, itemsep=2pt, topsep=1pt]
  \item 若将上述微扰项修改为 $H'' = \lambda (x^2 + y^2) z$，基态的一级能量修正 $E_0^{(1)}$ 是否为零？为什么？
  \item 曲线积分 $\oint_L \frac{-y dx + x dy}{x^2+y^2}$ 中，虽然满足 $\frac{\partial Q}{\partial x} - \frac{\partial P}{\partial y} = 0$，但当闭合曲线 $L$ 包含原点时积分却为 $2\pi$。请从微积分定理前提与场论物理旋度两个维度简述其根本原因。
\end{enumerate}
\end{problem}

% ==============================================================================
% PAGE 4: §3 CALCULATION TOPIC (计算专题)
% ==============================================================================
\clearpage

{\Large\sffamily\bfseries\color{navy} \S 3\quad Calculation Topic \textbar\ 计算专题}\hfill{\small\sffamily\color{mist}\itshape Pure Calculation \& Tricks}

\vspace{8pt}

\begin{problem}{Level 1 \textbar\ Basic Execution: Green's Theorem Standard Conversion}{Suggested time: 10\,min}
计算曲线积分：
\[
  I_1 = \oint_L (2xy - y^2)\,dx + (x^2 + 2xy)\,dy
\]
其中 $L$ 为由抛物线 $y = x^2$ 与直线 $y = x$ 围成平面闭区域的正向（逆时针）边界。
\end{problem}

\vspace{6pt}

\begin{problem}{Level 2 \textbar\ Advanced Tricks: Supplementary Line Method (补线法)}{Suggested time: 15\,min}
计算曲线积分：
\[
  I_2 = \int_L (e^x \sin y - y)\,dx + (e^x \cos y - 1)\,dy
\]
其中 $L$ 为上半圆周 $x^2 + y^2 = a^2 \, (y \ge 0)$，曲线方向取从点 $(a, 0)$ 到点 $(-a, 0)$。
\end{problem}

\vspace{6pt}

\begin{problem}{Level 3 \textbar\ Boss Traps: Singularity \& Puncture Hole Method (奇点挖洞法)}{Suggested time: 20\,min}
计算曲线积分：
\[
  I_3 = \oint_L \frac{x dy - y dx}{x^2 + y^2} + (x+y)\,dx + (x-y)\,dy
\]
其中 $L$ 为椭圆曲线 $\frac{x^2}{4} + \frac{y^2}{9} = 1$，绕行方向为逆时针方向。
\end{problem}

% ==============================================================================
% PAGE 5: §4 MODEL SOLUTIONS & COMMENTARY
% ==============================================================================
\clearpage

{\Large\sffamily\bfseries\color{navy} \S 4\quad Model Solutions \& Commentary}

\vspace{1pt}
{\small\color{mist}\itshape\hspace{1pt} Complete your independent attempt on scratch paper before consulting this page.}

\vspace{6pt}

\begin{solution}{M1 \textperiodcentered\ Model Solution \& Commentary}
\textbf{Strategy:}\enspace 被积函数拆项后转极坐标计算，或结合区域 $D$ 关于直线 $y = x$ 的轮换对称性化简。

\vspace{1.5mm}
\textbf{Derivation:}\enspace
区域 $D$ 为第一象限四分之一圆盘，变换为极坐标：$x = r\cos\theta, y = r\sin\theta$，定限为：$0 \le r \le 1, \, 0 \le \theta \le \frac{\pi}{2}$，面积微元 $dxdy = r\,dr\,d\theta$。

被积函数化简：
\[
  \frac{x^3 + x^2 y}{x^2 + y^2} = \frac{r^3(\cos^3\theta + \cos^2\theta\sin\theta)}{r^2} = r(\cos^3\theta + \cos^2\theta\sin\theta)
\]
原二重积分化为累次积分：
\[
  I = \int_0^{\pi/2} (\cos^3\theta + \cos^2\theta\sin\theta)\,d\theta \cdot \int_0^1 r \cdot r\,dr
\]
分步计算各部分：
\begin{enumerate}[leftmargin=1.5em, itemsep=1pt, topsep=1pt]
  \item \textbf{径向积分}：$\int_0^1 r^2\,dr = \left[\frac{1}{3}r^3\right]_0^1 = \frac{1}{3}$。
  \item \textbf{角度积分}：
    \begin{itemize}[leftmargin=1.2em, itemsep=1pt]
      \item 第一项由 Wallis 公式：$\int_0^{\pi/2} \cos^3\theta\,d\theta = \frac{2}{3}$；
      \item 第二项凑微分：$\int_0^{\pi/2} \cos^2\theta\sin\theta\,d\theta = \left[-\frac{1}{3}\cos^3\theta\right]_0^{\pi/2} = \frac{1}{3}$。
    \end{itemize}
    两项相加：$\int_0^{\pi/2} (\cos^3\theta + \cos^2\theta\sin\theta)\,d\theta = \frac{2}{3} + \frac{1}{3} = 1$。
\end{enumerate}
综合以上两部分乘积，得最终结果：
\[
  I = 1 \times \frac{1}{3} = \frac{1}{3}
\]

\begin{notebox}
\alerttag{Alternative Method}\enspace 亦可利用轮换对称性：$\iint_D \frac{x^3}{x^2+y^2}dxdy = \iint_D \frac{y^3}{x^2+y^2}dxdy = \frac{1}{2}\iint_D \frac{x^3+y^3}{x^2+y^2}dxdy$，进而结合极坐标化简。
\end{notebox}
\end{solution}

\vspace{4pt}

\begin{solution}{Q1 \textperiodcentered\ Model Solution \& Commentary (2015 UCAS 真题)}
\textbf{Physical Picture:}\enspace 微扰项为坐标三次奇宇称项，对角矩阵元必为零；二阶能量由升降算符筛选仅激发 $|010\rangle, |100\rangle, |210\rangle, |120\rangle$ 四个中间态。

\vspace{1.5mm}
\textbf{Derivation:}\enspace
(1) 宇称反演 $x \to -x, y \to -y$，微扰算符 $H'(-x,-y) = -H'(x,y)$ 为奇宇称算符。基态 $|000\rangle$ 概率密度为偶函数，故 $E_0^{(1)} = \langle 000|H'|000\rangle = 0$。

(2) 设 $C = \sqrt{\frac{\hbar}{2m\omega}}$，则 $x = C(a_x+a_x^\dagger), y = C(a_y+a_y^\dagger)$。
微扰作用于基态：
\[
  x^2 y |000\rangle = C^3 (|0\rangle_x + \sqrt{2}|2\rangle_x)|1\rangle_y |0\rangle_z = C^3 (|010\rangle + \sqrt{2}|210\rangle)
\]
同理 $x y^2 |000\rangle = C^3 (|100\rangle + \sqrt{2}|120\rangle)$。非零中间态仅 4 个：
\begin{itemize}[leftmargin=1.5em, itemsep=1pt, topsep=1pt]
  \item $|010\rangle, |100\rangle$：分母 $\Delta E = E_{000}^{(0)} - E_{010}^{(0)} = -\hbar\omega$，模方和为 $2 \lambda^2 C^6$；
  \item $|210\rangle, |120\rangle$：分母 $\Delta E = -3\hbar\omega$，模方和为 $2 \times 2\lambda^2 C^6 = 4\lambda^2 C^6$。
\end{itemize}
求和合并：
\[
  E_0^{(2)} = -\frac{2\lambda^2 C^6}{\hbar\omega} - \frac{4\lambda^2 C^6}{3\hbar\omega} = -\frac{10\lambda^2 C^6}{3\hbar\omega}
\]
代入 $C^6 = \left(\frac{\hbar}{2m\omega}\right)^3 = \frac{\hbar^3}{8m^3\omega^3}$，化简得：$E_0^{(2)} = -\frac{5\lambda^2 \hbar^2}{12 m^3 \omega^4}$。

(3) 因 $E_0^{(2)} < 0$，基态能量下移。物理图像：二阶微扰中基态处于能谱最底层，所有上方激发态均通过微扰“排斥”基态，将其能级向下推移。
\end{solution}

\vspace{4pt}

\begin{solution}{Calculation Topic \textperiodcentered\ Solutions}
\textbf{Level 1:}\enspace $P = 2xy - y^2, Q = x^2 + 2xy$。$\frac{\partial Q}{\partial x} - \frac{\partial P}{\partial y} = (2x+2y) - (2x-2y) = 4y$。
由格林公式：$I_1 = \iint_D 4y dxdy = \int_0^1 dx \int_{x^2}^x 4y dy = \int_0^1 2(x^2 - x^4)dx = 2(\frac{1}{3} - \frac{1}{5}) = \frac{4}{15}$。

\vspace{1.5mm}
\textbf{Level 2:}\enspace 添补线段 $L_1: y=0, x$ 从 $-a$ 到 $a$。闭曲线 $C = L + L_1$（逆时针）。
$\frac{\partial Q}{\partial x} - \frac{\partial P}{\partial y} = e^x\cos y - (e^x\cos y - 1) = 1$。
$\oint_C Pdx+Qdy = \iint_D 1 dxdy = \frac{1}{2}\pi a^2$。在 $L_1$ 上 $y=0, dy=0 \implies \int_{L_1} = 0$。
故 $I_2 = \oint_C - \int_{L_1} = \frac{1}{2}\pi a^2$。

\vspace{1.5mm}
\textbf{Level 3:}\enspace 拆为两部分：$I_{3A} = \oint_L \frac{xdy - ydx}{x^2+y^2}$，第二部分 $P_2 = x+y, Q_2 = x-y \implies \frac{\partial Q_2}{\partial x} - \frac{\partial P_2}{\partial y} = 1 - 1 = 0 \implies I_{3B} = 0$。
对于 $I_{3A}$，原点 $(0,0)$ 为奇点。在椭圆内挖小圆 $l_\varepsilon: x^2+y^2=\varepsilon^2$（顺时针），在复连通区域内无奇点，旋度为 0。由挖洞法：$I_{3A} = \oint_{l_\varepsilon \text{(逆时针)}} \frac{xdy-ydx}{x^2+y^2} = 2\pi$。最终结果 $I_3 = 2\pi + 0 = 2\pi$。
\end{solution}

% ==============================================================================
% PAGE 6: §5 DAILY DEBRIEF & REFLECTION
% ==============================================================================
\clearpage

{\Large\sffamily\bfseries\color{navy} \S 5\quad Daily Debrief \& Reflection}

\vspace{6pt}

\begin{debrief}{Post-Match Review (fill in before sleep, 5\,min)}
\renewcommand{\arraystretch}{1.2}
\sffamily\small
\begin{tabularx}{\linewidth}{@{}p{2.4cm}X@{}}
  \textbf{604 Calculus} & Time spent: \underline{\hspace{1.4cm}} min \quad\textbar\quad M1 status: \checkbox\ Independent \checkbox\ Used hints \checkbox\ Incomplete\\
  & Error type: \checkbox\ Knowledge gap \checkbox\ Modelling \checkbox\ Logic \checkbox\ Calculation \checkbox\ Time\\
  \midrule
  \textbf{811 QM} & Time spent: \underline{\hspace{1.4cm}} min \quad\textbar\quad Q1 status: \checkbox\ Independent \checkbox\ Used hints \checkbox\ Incomplete\\
  & Error type: \checkbox\ Knowledge gap \checkbox\ Physical model \checkbox\ Operator algebra \checkbox\ Calculation \checkbox\ Time\\
  \midrule
  \textbf{Calc Topic} & Time spent: \underline{\hspace{1.4cm}} min \quad\textbar\quad Level 1 \checkbox\ \textbar\ Level 2 \checkbox\ \textbar\ Level 3 \checkbox\ \textbar\ All correct \checkbox\\
  \midrule
  \textbf{Planning} & \textbf{Spaced retrieval in 3 days:} \underline{\hspace{8cm}}\\
  & \textbf{Today's assessment:} \checkbox\ Peak performance \checkbox\ Steady progress \checkbox\ Minimum defence mode\\
\end{tabularx}
\end{debrief}

\vspace{8pt}

\begin{tcolorbox}[
  enhanced, colback=white, colframe=border, boxrule=0.6pt, arc=2mm,
  left=4mm, right=4mm, top=3.5mm, bottom=4mm,
  fonttitle=\sffamily\bfseries\color{slate},
  title={\raisebox{0pt}{\color{forest}\rule{3pt}{10pt}}\hspace{4pt}Key Derivation Insights \& Traps \textbar\ 突破口与避坑沉淀},
  colbacktitle=paper, coltitle=slate,
  titlerule=0pt, toptitle=2mm, bottomtitle=2mm
]
\sffamily\small
\textbf{Core Breakthrough (今日核心突破口与关键一步)}:\\[30pt]
\rule{\linewidth}{0.3pt}\\[12pt]
\textbf{Fatal Trap \& Common Error (致命陷阱与避坑警示)}:\\[30pt]
\rule{\linewidth}{0.3pt}\\[12pt]
\textbf{Day +3 Action Item (3 天后间隔重做提取的具体题号与断点)}:\\[25pt]
\rule{\linewidth}{0.3pt}
\end{tcolorbox}

\end{document}
